%-------------------------
% Resume in Latex
% Author : Jake Gutierrez
% Based off of: https://github.com/sb2nov/resume
% License : MIT
%------------------------

\documentclass[letterpaper,11pt]{article}

\usepackage{latexsym}
\usepackage[empty]{fullpage}
\usepackage{titlesec}
\usepackage{marvosym}
\usepackage[usenames,dvipsnames]{color}
\usepackage{verbatim}
\usepackage{enumitem}
\usepackage[hidelinks]{hyperref}
\usepackage{fancyhdr}
\usepackage[english]{babel}
\usepackage{tabularx}
\input{glyphtounicode}


%----------FONT OPTIONS----------
% sans-serif
% \usepackage[sfdefault]{FiraSans}
% \usepackage[sfdefault]{roboto}
% \usepackage[sfdefault]{noto-sans}
% \usepackage[default]{sourcesanspro}

% serif
% \usepackage{CormorantGaramond}
% \usepackage{charter}
% \usepackage{times}

\usepackage[cmintegrals,cmbraces]{newtxmath}
\usepackage{ebgaramond-maths}
\usepackage[T1]{fontenc}

\pagestyle{fancy}
\fancyhf{} % clear all header and footer fields
\fancyfoot{}
\renewcommand{\headrulewidth}{0pt}
\renewcommand{\footrulewidth}{0pt}

% Adjust margins
\addtolength{\oddsidemargin}{-0.5in}
\addtolength{\evensidemargin}{-0.5in}
\addtolength{\textwidth}{1in}
\addtolength{\topmargin}{-.5in}
\addtolength{\textheight}{1.0in}

\urlstyle{same}

\raggedbottom
\raggedright
\setlength{\tabcolsep}{0in}

% Sections formatting
\titleformat{\section}{
  \vspace{-4pt}\scshape\raggedright\large
}{}{0em}{}[\color{black}\titlerule \vspace{-5pt}]

% Ensure that generate pdf is machine readable/ATS parsable
\pdfgentounicode=1

%-------------------------
% Custom commands
\newcommand{\resumeItem}[1]{
  \item\small{
    {#1 \vspace{-2pt}}
  }
}

\newcommand{\resumeSubheading}[4]{
  \vspace{-2pt}\item
    \begin{tabular*}{0.97\textwidth}[t]{l@{\extracolsep{\fill}}r}
      \textbf{#1} & #2 \\
      \textit{\small#3} & \textit{\small #4} \\
    \end{tabular*}\vspace{-7pt}
}

\newcommand{\resumeSubSubheading}[2]{
    \item
    \begin{tabular*}{0.97\textwidth}{l@{\extracolsep{\fill}}r}
      \textit{\small#1} & \textit{\small #2} \\
    \end{tabular*}\vspace{-7pt}
}

\newcommand{\resumeProjectHeading}[2]{
    \item
    \begin{tabular*}{0.97\textwidth}{l@{\extracolsep{\fill}}r}
      \small#1 & #2 \\
    \end{tabular*}\vspace{-7pt}
}

\newcommand{\resumeSubItem}[1]{\resumeItem{#1}\vspace{-4pt}}

\renewcommand\labelitemii{$\vcenter{\hbox{\tiny$\bullet$}}$}

\newcommand{\resumeSubHeadingListStart}{\begin{itemize}[leftmargin=0.15in, label={}]}
\newcommand{\resumeSubHeadingListEnd}{\end{itemize}}
\newcommand{\resumeItemListStart}{\begin{itemize}}
\newcommand{\resumeItemListEnd}{\end{itemize}\vspace{-5pt}}

%-------------------------------------------
%%%%%%  RESUME STARTS HERE  %%%%%%%%%%%%%%%%%%%%%%%%%%%%


\begin{document}

%----------HEADING----------
\begin{nohyperpage}
\begin{tabularx}{\textwidth}{@{}X r@{}}
  \begin{tabular}[t]{@{}l@{}}
    \textbf{\Huge Emre Kerman} \\ \vspace{2pt}
    {\large\textit{Bilgisayar Mühendisi}} \\ \vspace{1pt}
    \small İzmir, Türkiye
  \end{tabular} & 
  \begin{tabular}[t]{r@{}}
    \small +90 541 264 0344 \\
    \href{mailto:emrekerman02@gmail.com}{\underline{emrekerman02@gmail.com}} \\
    \href{https://github.com/emrekerman}{\underline{github.com/emrekerman}} \\
    \href{https://linkedin.com/in/kermanemre}{\underline{linkedin.com/in/kermanemre}}
  \end{tabular}
\end{tabularx}
\end{nohyperpage}
\vspace{-10pt}


%-------------------------------------------
\section{Özet}
Dokuz Eylül Üniversitesi'nden yeni mezun olmuş bir Bilgisayar Mühendisiyim. Yazılım geliştirme, makine öğrenmesi, gömülü sistemler ve full-stack web geliştirme alanlarında deneyim sahibiyim. IoT ve computer vision sistemleri, frontend ve backend uygulamaları geliştirme; Linux ve açık kaynak teknolojilerle çalışma konusunda deneyimliyim.


%-----------EĞİTİM-----------
\section{Eğitim}
  \resumeSubHeadingListStart
    \resumeSubheading
      {Dokuz Eylül Üniversitesi}{İzmir, Türkiye}
      {Bilgisayar Mühendisliği Lisans Derecesi}{2021 -- 2025}
    \resumeSubheading
      {Yunus Emre Anadolu Lisesi}{İzmir, Türkiye}
      {Lise Diploması}{2018 -- 2021}
  \resumeSubHeadingListEnd


%-----------DENEYİM-----------
\section{Deneyim}
  \resumeSubHeadingListStart

    \resumeSubheading
      {Yazılım Mühendisliği Stajyeri}{Haziran 2025 -- Eylül 2025}
      {Indas Teknoloji Ticaret AŞ.}{İzmir, Türkiye}
      \resumeItemListStart
        \resumeItem{ESP32-CAM node'ları ve OCR (PaddleOCR/Tesseract) kullanarak otomatik araç takibi, ücretlendirme ve bakiye yönetimi sağlayan akıllı otopark platformu geliştirdim.}
        \resumeItem{Quantized FOMO modelleri, C++/Arduino ve Edge Impulse kullanarak ESP32-CAM üzerinde çalışan edge-first bir Kişisel Koruyucu Donanım (KKD) ve insan tespit sistemi geliştirdim; 120--250 ms inference gecikmesi, 18 FPS canlı yayın ve 5--6 FPS cihaz üzerinde inference performansına ulaştım.}
        \resumeItem{Gerçek zamanlı bounding box gösterimi, tespit kayıtları ve bakiye panelleri için Next.js ve React kullanarak web arayüzleri geliştirdim.}
      \resumeItemListEnd

    \resumeSubheading
      {Yazılım Mühendisliği Stajyeri}{Kasım 2025 -- Aralık 2025}
      {Dokuz Eylül Üniversitesi Bilgisayar Mühendisliği Bölümü}{İzmir, Türkiye}
      \resumeItemListStart
        \resumeItem{"Restock" adlı kiosk ve mobil cihazlar arasında gerçek zamanlı oturum devri sağlayan self-checkout platformunu geliştirdim.}
        %\resumeItem{Fiziksel mağaza kiosklarını mobil cihazlarla entegre ederek gerçek zamanlı oturum devri sağlayan "Restock" adlı self-checkout platformunu geliştirdim.}
        \resumeItem{Drizzle ORM ve SQLite kullanarak atomic veritabanı işlemleri uyguladım ve barkod tarama için özel React hook'ları geliştirdim.}
        \resumeItem{Next.js Middleware üzerinden JWT tabanlı rol bazlı yetkilendirme yapılandırdım ve yönetimsel izleme panelleri geliştirdim.}
      \resumeItemListEnd

  \resumeSubHeadingListEnd
  
%-----------PROJELER-----------
\section{Projeler}
    \resumeSubHeadingListStart
      \resumeProjectHeading
          {\textbf{Early Frost Detection} $|$ \emph{Go, Python, C, TensorFlow, Keras, NumPy, ESP32}}{2024 -- 2025}
          \resumeItemListStart
            \resumeItem{Bir Arçelik kullanım senaryosu için buzdolabı soğutma kanatçıklarında buz oluşumunu tahmin eden makine öğrenmesi tabanlı bir IoT çözümü geliştirdim.}
            \resumeItem{ESP32 mikrodenetleyiciler ve DHT22 sensörleri kullanarak bir buzdolabı prototipi yaptım; Wi-Fi üzerinden çevresel veri topladım.}
            \resumeItem{Toplanan veriler üzerinde veri ön işleme, özellik mühendisliği ve makine öğrenmesi modeli eğitimi gerçekleştirdim.}
            \resumeItem{Go kullanarak gerçek zamanlı bir web sunucusu ve frontend izleme platformu geliştirdim.}
          \resumeItemListEnd

      \resumeProjectHeading
          {\textbf{Hand Gesture Detection} $|$ \emph{Python, TensorFlow, Keras, NumPy, OpenCV}}{2024}
          \resumeItemListStart
            \resumeItem{Web kamerası üzerinden el hareketlerini sınıflandırmak için hafif bir CNN ve ResNet benzeri bir mimari olmak üzere iki computer vision modeli geliştirdim.}
            \resumeItem{Model eğitimi için özel bir el hareketi veri seti oluşturdum ve ön işleme tabi tuttum.}
            \resumeItem{Veri kaybını önlemek için Model Checkpoint'lerini kullanan verimli bir eğitim pipeline'ı geliştirdim.}
            \resumeItem{Gerçek zamanlı inference için canlı web kamerası girdisini kullanıcı deneyimi iyileştirmeleriyle entegre ettim.}
          \resumeItemListEnd

      \resumeProjectHeading
          {\textbf{Game Platform \& Storefront} $|$ \emph{C\#, .NET Core, PostgreSQL}}{2023}
          \resumeItemListStart
            \resumeItem{Dijital oyun dağıtım platformu için veritabanı mimarisini tasarlayıp geliştirdim.}
            \resumeItem{Kullanıcı kimlik doğrulama, oyun yayınlama ve satın alma işlemlerini destekleyen backend uygulama mantığını geliştirdim.}
          \resumeItemListEnd
    \resumeSubHeadingListEnd

 %-----------TECHNICAL SKILLS-----------
\section{Teknik Yetkinlikler}
 \begin{itemize}[leftmargin=0.15in, label={}]
    \small{\item{
     \textbf{Programlama Dilleri}{: TypeScript, JavaScript, Python, C/C++, Go, C\#, SQL, Rust} \\
     \textbf{Web}{: Next.js, React, React Native, Astro, Tailwind CSS, REST APIs} \\
     \textbf{Backend ve Veritabanları}{: PostgreSQL, SQLite, Drizzle ORM, .NET Core, OAuth/JWT, Swagger/OpenAPI} \\
     \textbf{Yapay Zeka, Görüntü İşleme ve Edge}{: TensorFlow, Keras, TFLite, OpenCV, Edge Impulse, ESP32/ESP32-CAM, PaddleOCR} \\
     \textbf{Sistemler ve DevOps}{: Linux, Docker, Kubernetes, Git, Nix/NixOS} \\
     \textbf{Diller}{: Türkçe (Ana Dil), İngilizce (C1), Almanca (A1)}
    }}
 \end{itemize}

%-------------------------------------------
\end{document}
